\section{Order and Progress}

\begin{quotex}
The feebleness of our intelligence, and the shortness of individual life compared with the slowness of social development, keep our imagination, especially as regards political ideas, which are very complex, entirely dependent on the milieu in which we actually live. Even the most declared utopians, who believe they have emancipated themselves from every condition of reality, undergo, all unconsciously this necessity, and faithfully reflect in their dreams the contemporary social state. A fortiori, the conception of a political system radically different from the one that surrounds us is quite beyond the bounds of our intelligence.
\end{quotex}

The movement of thought following \textbf{Francis Bacon} opposed Progress to the Order of the Medieval period. In the aftermath of the French Revolution, another world figure emerged, \textbf{Auguste Comte}, who tried to think through the problem of Progress and Order in a deeper way. Unwilling to forego the benefits of science based on the methodology instaurated by Bacon, neither was he willing ignore the natural and cosmic Order. He concluded that real progress is just the unfolding in time of a primeval order much as the oak tree is the fulfillment of the acorn. In his own words:

\begin{quotationx}
All previous philosophies had regarded Order as stationary, a conception which rendered it wholly inapplicable to modern politics. But Positivism, by rejecting the absolute, and yet not introducing the arbitrary, represents Order in a totally new light, and adapts it to our progressive civilization. It places it on the firmest possible foundation, that is, on the doctrine of the invariability of the laws of nature, which defends it against all danger from subjective chimeras. The Positivist regards artificial Order in Social phenomena, as in all others, as resting necessarily upon the Order of nature, in other words, upon the whole series of natural laws.

But Order has to be reconciled with Progress; and here Positivism is still more obviously without a rival. Necessary as the reconciliation is, no other system has even attempted it. But the facility with which we are now enabled, by the encyclopedic scale, to pass from the simplest mathematical phenomena to the most complicated phenomena of political life, leads at once to a solution of the problem... . Order is the condition of all Progress; Progress is always the object of Order. Or, to penetrate the question still more deeply, Progress may be regarded simply as the development of Order; for the order of nature necessarily contains within itself the germ of all possible progress. The rational view of human affairs is to look on all their changes, not as new Creations, but as new Evolutions. And we find this principle fully borne out in history. ... Progress then is in its essence identical with Order, and may be looked upon as Order made manifest.
\end{quotationx}


He called his philosophical conception \textbf{Positivism} since it is based solely on positive facts. From a rather Eurocentric perspective, he postulated that society pass through three stages. The first is the theological (subdivided into polytheism and monotheism), then the ideological (which he calls the metaphysical). The first stage, which comprises Greco-Roman Antiquity and the Middle Ages, attributes everything to divine beings. The ideological stage, making up the Renaissance and Enlightenment, proposes abstract notions as the explanatory principle. (Gornahoor reserves the term metaphysics for something so different.) Positivism claims to be based on nothing but empirical facts, which may include inner states. To his credit, he avoids ideologies such as materialism, determinism, reductionism, and the like, which most philosophers tend to conflate with Positivism.

He organizes the sciences in an ascending sequence of a more general and complex order: mathematics, astronomy, physics, chemistry, biology, sociology (including political science), and eventually morality. Thus he comes close to a complete understanding of ideas (See Ideas and Occult History\footnote{\url{https://gornahoor.net/?p=3451}}), though he never takes the next step. He claimed to have developed the new science of sociology, which he hoped would be the key to true Progress based on Order. We pointed out the importance of a mathematical foundation in Maths and Politics\footnote{\url{https://www.gornahoor.net/?p=37}}. Since Comte's time, there have been significant improvements in mathematics such as game theory, actuarial science, chaos theory which have untapped significance for understanding the behavior of crowds and its ramifications. Nevertheless, the field is still mired in ideology, which retards sociology and political science at a rather primitive level. Comte proposes four degrees of Progress, in order of increasing generality and complexity in the phenomena:

\begin{enumerate}


\item \textbf{Material Improvement} (economic activity, industry)


\item \textbf{Physical Improvement} (health and longevity)


\item \textbf{Intellectual Improvement} (knowledge, the True)


\item \textbf{Moral Improvement} (will, the Good)

\end{enumerate}


Material improvement is straight-forward, keeping in mind he means for the common good and not for the good of a small oligarchy. The Physical Improvement of the human race seems reasonable when restricted to the health and longevity of the individual. However, from the perspective of the longer timeframe of social development, this is a project that needs to be thought out more carefully. In a closed system\footnote{\url{https://en.wikipedia.org/wiki/Systems\_theory}}, such as a social group, what may seem to be good from the local perspective may be deleterious to the system as a whole.

Intellectual and Moral improvement have been a frequent topic, and our successes at this, outside of the developments of applied science, have been limited. Comte noted that sociology must be guided by the three fundamental elements of human nature: Feeling, Thought, and Action (or Will)\footnote{\url{https://www.gornahoor.net/?p=2683}}. The overwhelming mass of humanity is driven by the life of Feeling. Men want to do what feels good, gives them social status, or puffs up their ego. The intellectual and moral life is opposed to that, and is a product of spiritual development\footnote{\url{https://www.meditationsonthetarot.com/the-vow-of-poverty}}. The commitment to the rational or the good, whenever it goes against the life of Feeling, will reveal its weakness. Comte writes:

\begin{quotex}
Owing to the mental and moral anarchy in which we live, systematic efforts to gain the higher degrees of Progress are as yet impossible; and this explains, though it does not justify, the exaggerated importance attributed nowadays to material improvements. But the only kinds of improvement really characteristic of Humanity are those which concern our own nature... 
\end{quotex}


Nevertheless, he held out hope to preserve the best in Order and Progress:

Positivism, then, realizes the highest aspirations of medieval Catholicism ... it combines the opposite merits of the Catholic and Revolutionary spirit, and by so doing supersedes them both. Theology and Metaphysics Ideology may now disappear without danger, because the service which each of them rendered is now harmonized with that of the other, and will be performed more perfectly.

Unfortunately, in ascending the ladder of general systems, Comte never reached the ultimate and all-inclusive idea, the Logos. Here, we would crown Comte's system with Metaphysics, in its Traditional meaning, with its own Metaphysical Positivism\footnote{\url{https://www.gornahoor.net/?p=1189}}, as \textbf{Julius Evola} proposed. Perhaps, then, men would arise who have the intellectual power and strength of Will to bring about Progress based on Order. Then Comte's program may be revived.

\begin{quotex}
Positivists will always acknowledge the close relation between their own system and the memorable effort of medieval Catholicism. In offering for the acceptance of Humanity a new organization of life, we would not dissociate it with all that has gone before. On the contrary, it is our boast that we are but proposing for her maturity the accomplishment of the noble effort of her youth, an effort made when intellectual and social conditions precluded the possibility of success.
\end{quotex}



\flrightit{Posted on 2011-11-29 by Cologero}

